\chapter{羰基：亲核加成、缩醛与保护}\label{ch:b1:acetals-protection}

将葡萄糖溶于水后，几乎
没有分子以教科书中画出的开链醛形式存在：分子会弯曲，其自身的羟基
加到醛基上，闭合成一个六元环。醇对羰基的同样加成，
若再重复一次并以酸为催化剂，便得到缩醛。
缩醛对碱以及最活泼的亲核试剂均稳定，
这种稳定性是一种工具：一个会被格氏试剂破坏的醛，
可先掩蔽为缩醛，穿过整个反应，最后再复原。
本章研究羰基作为亲电试剂的性质，
以及水与醇的加成，
还有保护策略。

\begin{recall}
羰基化合物的格氏加成（见\cref{ch:b1:organomagnesium}）；
反应商 $Q$ 以及 $Q \neq K$ 时体系的演化
（见\cref{ch:b1:extent-q-and-k}）；弯箭头、亲核试剂与
亲电试剂（见\cref{ch:b1:electronic-effects}）。第 1 册（十二
年级）介绍过合成过程中保护基团的想法。
\end{recall}

\section{羰基作为亲电试剂}

C=O 键向氧一侧极化（见\cref{ch:b1:periodicity}），且
其共振结构 \ce{C+-O-} 使碳带满正电荷：
碳具亲电性，亲核试剂从垂直于
羰基平面的方向进攻。强亲核试剂（格氏试剂、氢化物、
氰根）可直接加成。弱亲核试剂——水、醇——则需要帮助。

\begin{definition}[亲电活化]\label{def:b1:acetals-protection:activation}
\emph{亲电活化}\index{亲电活化}是指羰基在
氧原子上发生质子化（或与路易斯酸配位）。
生成的阳离子 \ce{C=OH+} 的正电荷可经共振结构 \ce{C+-OH}
由碳分享，因而即使是弱
亲核试剂也能进攻它。
\end{definition}

\begin{omfigure}
\setchemfig{atom sep=2.1em}
\schemestart
\chemfig{C(-[:150]R)(-[:210]H)=O}
\arrow{->[\ce{H+}]}
\chemfig{C(-[:150]R)(-[:210]H)=O^{+}-H}
\arrow{<->}
\chemfig{C^{+}(-[:150]R)(-[:210]H)-O-H}
\schemestop

{\small 酸催化下醛的活化：氧原子质子化
生成阳离子，其共振结构使电荷
位于碳原子上。}
\end{omfigure}

\section{水合物与半缩醛}

水可与醛和酮发生可逆加成，\ce{R2C=O + H2O <=>
R2C(OH)2}；水合物对酮通常占少数，对
甲醛则可能占主要部分。醇的加成方式相同。

\begin{definition}[半缩醛与缩醛]\label{def:b1:acetals-protection:hemiacetal}
\emph{半缩醛}\index{半缩醛}是指一个碳原子
同时带有一个 OH 基和一个 OR 基的化合物，由醇对羰基加成而成。
\emph{缩醛}\index{缩醛}是指一个碳原子带有两个 OR 基的化合物，
由半缩醛与第二分子醇生成，
同时失去一分子水。
\end{definition}

开链半缩醛与醛及醇的平衡中通常占少数；
当羟基与 C=O 位于同一分子中，
且能闭合成五元环或六元环时，
环状半缩醛占主导。葡萄糖即是如此，
较简单的 5-羟基戊醛也是如此。

\begin{omfigure}
\setchemfig{atom sep=2em}
\schemestart
\chemfig{HO-[:30]-[:-30]-[:30]-[:-30]-[:30]C(=[:90]O)-[:-30]H}
\arrow{<=>}
\chemfig{*6(O-(-[:-30]OH)-----)}
\schemestop

\medskip

{\small 5-羟基戊醛及其环状半缩醛，为一个含氧的六元环。
碳 5 的羟基已加到碳 1 的醛基上；
与葡萄糖相同，环式占主导。}
\end{omfigure}

\section{缩醛}

\begin{proposition}[缩醛生成的机理]\label{prop:b1:acetals-protection:acetal-mechanism}
在酸性条件下，醛或酮与醇先生成半缩醛，
再生成缩醛，步骤如下：C=O 质子化；醇加成；
失去质子（得半缩醛）；羟基质子化；失去水，
生成由剩余 OR 基稳定的碳正离子（即氧正离子）；
加入第二分子醇；失去质子。每一步均
可逆。
\end{proposition}

\begin{proof}
每一步或为酸碱质子转移，或为亲核试剂对亲电碳的加成
（或其逆过程），如
\cref{ch:b1:electronic-effects}中所述。失去水之所以可能，是因为所生成的
碳正离子 $\mathrm{R_2C^+{-}OR}$ 具有共振结构 $\mathrm{R_2C{=}O^+R}$，
其中每个原子均满足八隅体（见\cref{prop:b1:electronic-effects:carbocation-order}）。
酸在质子化步骤中被消耗，
在去质子化步骤中再生，因此酸是催化剂。
\end{proof}

\begin{omfigure}
\setchemfig{atom sep=1.9em}
\schemestart
\chemfig{C(-[:150]R)(-[:210]H)=O}
\arrow{->[\ce{H+}][]}
\chemfig{C^{+}(-[:150]R)(-[:210]H)-OH}
\arrow{->[$\mathrm{R'OH}$][\ce{-H+}]}
\chemfig{C(-[:150]R)(-[:210]H)(-[2]OR')-OH}
\schemestop

\medskip
\schemestart
\chemfig{C(-[:150]R)(-[:210]H)(-[2]OR')-OH}
\arrow{->[\ce{H+}][\ce{-H2O}]}
\chemfig{C^{+}(-[:150]R)(-[:210]H)-OR'}
\arrow{->[$\mathrm{R'OH}$][\ce{-H+}]}
\chemfig{C(-[:150]R)(-[:210]H)(-[2]OR')-OR'}
\schemestop
\par\vspace{6pt}

{\small 缩醛的生成分两步：先生成半缩醛（活化、
加入醇、失去一个质子），再生成缩醛（羟基质子化、
失水生成稳定的阳离子、加入第二分子醇、
失去一个质子）。}
\end{omfigure}

总反应 $\mathrm{RCHO} + 2\,\mathrm{R'OH} \rightleftharpoons \mathrm{RCH(OR')_2} + \ce{H2O}$ 的平衡常数
较小，且水为生成物。

\begin{proposition}[推动缩醛化反应]\label{prop:b1:acetals-protection:dean-stark}
若不断移去生成的水，缩醛化将持续进行，
直至羰基化合物（或醇）耗尽。
\end{proposition}

\begin{proof}
反应商 $Q = [\text{缩醛}][\ce{H2O}]/([\text{醛}]
[\mathrm{R'OH}]^2)$ 只要水被不断移去，便保持小于 $K$：由
\cref{ch:b1:extent-q-and-k}可知，反应持续正向进行，无论
$K$ 的值如何。若用二元醇（乙烷-1,2-二醇），缩醛为环状，
平衡也更有利，一分子二元醇替代两分子
醇。
\end{proof}

\begin{omfigure}
\begin{tikzpicture}[font=\small]
  \pic at (0,0) {heatingmantle};
  \pic at (0,0) {rbflask={0.4}{omDef!14}};
  \pic (d) at (0,3.0) {deanstark={0.35}{omDef!35}{orange!15}};
  \pic at (0,6.4) {condenser};
  \draw[->] (2.4,4.0) node[right, align=left] {收集器：水（下层）\\不再回流；甲苯 \\溢回烧瓶} -- (0.8,3.8);
  \draw[->] (-2.2,1.0) node[left, align=right] {醛、二元醇、甲苯、\\酸催化剂、沸腾} -- (-0.7,0.9);
  \draw[->] (1.8,8.0) node[right] {冷凝器} -- (0.4,7.8);
\end{tikzpicture}

{\small 迪安–斯塔克分水器。甲苯与水的蒸气冷凝并
落入带刻度的收集器中；水密度较大，且与甲苯不互溶，
下沉并留存，甲苯则溢回烧瓶。
收集到的水的体积即可度量反应的进程。}
\end{omfigure}

\begin{proposition}[缩醛的稳定性]\label{prop:b1:acetals-protection:acetal-stability}
缩醛在中性和碱性介质中稳定，对亲核试剂、
格氏试剂和氢化物试剂同样稳定；水溶液中的酸使其水解，
重新生成羰基化合物和醇。
\end{proposition}

\begin{proof}
缩醛碳上不带有可被碱或亲核试剂
驱除的离去基团：\ce{RO-} 为强碱，不会离去，并且
已不存在可被进攻的 C=O。在水溶液中，上述机理逆向进行：
过量的水使逆反应的 $Q < K$。
\end{proof}

\section{保护与脱保护}

\begin{definition}[保护基]\label{def:b1:acetals-protection:protecting-group}
称\emph{保护基}\index{保护基}为引入分子中的基团，
用以暂时掩蔽某一官能团，使其
在针对分子其他部分的步骤中不发生反应。其引入即为
\emph{保护}\index{保护}，其移除即为
\emph{脱保护}\index{脱保护}。
\end{definition}

\begin{method}[保护、反应、脱保护]\label{met:b1:acetals-protection:protect}
\begin{enumerate}
  \item 确定在所计划的步骤中会发生反应，或会
        破坏试剂的官能团。
  \item 选择在该步骤中稳定，且在
        分子其余部分可耐受的条件下可移除的保护基（对醛或
        酮面对碱、亲核试剂或氢化物时，用环状缩醛）。
  \item 保护；进行拟定步骤；脱保护。
  \item 需计入代价：多出两步，每步各有其产率。
\end{enumerate}
\end{method}

\begin{example}[碍事的酮]\label{ex:b1:acetals-protection:ketoester}
欲向一分子醛基加入格氏试剂，而该分子
还带有一个酮基，无简单的缩醛方案可选（醛比
酮更易生成缩醛，故会误保护错误的基团）；
须改变步骤顺序或所用试剂。保护是一种
策略，而非条件反射。
\end{example}

\section{练习}

\begin{exercise}[$\star$]\label{exo:b1:acetals-protection:1}
指出下列各式中属于半缩醛、缩醛、水合物、醚或醇的碳原子：
\ce{CH3CH(OH)OCH3}, \ce{CH3CH(OCH3)2}, \ce{CH3CH(OH)2}, \ce{CH3OCH3},
\ce{(CH3)2C(OCH2CH3)2}.
\end{exercise}

\begin{exercise}[$\star$]\label{exo:b1:acetals-protection:2}
写出丙醛与甲醇生成缩醛的总反应式，
以及丙酮与乙烷-1,2-二醇生成环状缩醛的总反应式。
\end{exercise}

\begin{exercise}[$\star$]\label{exo:b1:acetals-protection:3}
下列试剂中，哪些不会使缩醛发生变化：氢氧化钠
溶液、甲基溴化镁、稀硫酸水溶液、硼氢化
钠？
\end{exercise}

\begin{exercise}[$\star$]\label{exo:b1:acetals-protection:4}
画出由 4-羟基丁醛生成的环状半缩醛。
该环的大小是多少？
\end{exercise}

\begin{exercise}[$\star\star$]\label{exo:b1:acetals-protection:5}
写出乙醛二甲基缩醛的酸催化生成机理，
并用弯箭头表示每一步，
说明酸被再生。
\end{exercise}

\begin{exercise}[$\star\star$]\label{exo:b1:acetals-protection:6}
写出 2,2-二甲氧基丙烷酸性水解的机理，
并说明为什么要用大量的水。
\end{exercise}

\begin{exercise}[$\star\star$]\label{exo:b1:acetals-protection:7}
对 \qty{0.250}{mol} 某种酮进行缩醛化反应，
并用迪安–斯塔克分水器跟踪。完全转化时应收集到
多大体积的水（密度约 \qty{1.0}{g/mL}）？一小时后，
已收集 \qty{3.2}{mL}，转化率是多少？
\end{exercise}

\begin{exercise}[$\star\star$]\label{exo:b1:acetals-protection:8}
对同一羰基化合物，说明为何与乙烷-1,2-二醇
生成的环状缩醛比与两分子甲醇生成的缩醛
更完全地形成。
\end{exercise}

\begin{exercise}[$\star\star$]\label{exo:b1:acetals-protection:9}
为什么在用格氏试剂处理被保护的化合物之前，
必须先除去（或中和）酸催化剂？
\end{exercise}

\begin{exercise}[$\star\star\star$]\label{exo:b1:acetals-protection:10}
提出 5-羟基-5-甲基己-2-酮的合成路线，
\ce{(CH3)2C(OH)CH2CH2COCH3}，以 4-氯丁-2-酮和丙酮为原料，
并利用保护基；逐步写出每步的试剂及其作用。
\end{exercise}

\begin{exercise}[$\star\star\star$]\label{exo:b1:acetals-protection:11}
用反应商说明，过量的甲醇同样能推动醛的缩醛化，
并将这一方法与
除水方法作比较。
\end{exercise}

\begin{exercise}[$\star\star\star$]\label{exo:b1:acetals-protection:12}
说明葡萄糖的环状半缩醛在酸性甲醇中为何会发生反应，
生成甲基缩醛（即甲基糖苷），而葡萄糖的其他羟基
却不会转化为醚。
\end{exercise}

\section{问题：带有醛基的格氏试剂}

\begin{problem}[{周末问题——为什么 4-溴丁醛不能生成格氏试剂，其以环状缩醛形式
加以保护，借助迪安–斯塔克分水器，
格氏试剂与丙酮的加成、脱保护，以及
所收集水的体积}]
\label{pb:b1:acetals-protection:1}
一位化学家需要 5-羟基-5-甲基己醛 \ce{(CH3)2C(OH)CH2CH2CH2CHO}，
计划由 4-溴丁醛 \ce{BrCH2CH2CH2CHO} 和丙酮来制备它。
她从 \qty{15.09}{g} 的 4-溴丁醛出发。摩尔质量
（\unit{g/mol}）：\ce{C4H7BrO} 150.9，\ce{C2H6O2} 62.0，\ce{C6H11BrO2} 194.9，
\ce{C3H6O} 58.0，\ce{C9H18O3} 174.0，\ce{C7H14O2} 130.0，\ce{H2O} 18.0；
水的密度为 \qty{0.997}{g/mL}。
% ledger: rho:water

\medskip
\noindent\textbf{第一部分——问题。}
\begin{enumerate}
  \item 写出 4-溴丁醛所对应的格氏试剂。
  \item 为什么该试剂不能存在？写出使其消失的
        反应。
  \item 需要保护哪一官能团，防范什么？
  \item 为什么适合用缩醛？
  \item 为什么优先选用环状缩醛（与乙烷-1,2-二醇生成）？
  \item 计算 4-溴丁醛的物质的量。
  \item 若过量 \qty{20}{\percent}，求所需二醇的质量。
\end{enumerate}

\medskip
\noindent\textbf{第二部分——保护。}
\begin{enumerate}[resume]
  \item 写出与乙烷-1,2-二醇生成缩醛的方程式。
  \item 酸催化剂（4-甲基苯磺酸）
        的作用是什么？
  \item 描述迪安–斯塔克分水器的工作过程。
  \item 用 $Q$ 和 $K$ 说明除去水为何能推动反应进行。
  \item 为什么回流到烧瓶中的是甲苯，而不是水？
  \item 化学家如何判断反应已经完全？
  \item 计算完全转化时应得到的被保护溴化物的质量。
\end{enumerate}

\medskip
\noindent\textbf{第三部分——格氏反应步骤。}
\begin{enumerate}[resume]
  \item 写出由被保护溴化物生成格氏试剂的
        反应。
  \item 写出其与丙酮的加成反应及生成的醇盐。
  \item 为什么该步骤的后处理必须是弱碱性或中性的（
        例如氯化铵水溶液），而不能是酸性的？
  \item 计算该步骤产率为 \qty{70}{\percent} 时
        被保护醇的质量。
\end{enumerate}

\medskip
\noindent\textbf{第四部分——脱保护。}
\begin{enumerate}[resume]
  \item 写出在酸性水溶液中脱保护的反应。
  \item 为什么叔醇在温和的酸性水溶液中能保持不变？
  \item 产物可通过其羟基与
        醛基加成闭合成六元环。画出该环状半缩醛。
  \item 计算脱保护产率为 \qty{90}{\percent} 时
        所得羟基醛的质量。
  \item 化学家从 4-溴丁醛出发，总产率为多少？
  \item 计算完全转化时，保护步骤中生成
        的水的质量。
  \item 说明完全转化时，在迪安–斯塔克分水器
        中收集到的水的体积。
\end{enumerate}
\end{problem}
